\documentclass[12pt, a4paper]{article}

\usepackage[utf8]{inputenc}
\usepackage[T1]{fontenc}
\usepackage[french]{babel}
\usepackage[mathscr]{euscript}
\let\euscr\mathscr \let\mathscr\relax
\usepackage[scr]{rsfso}
\usepackage{amssymb,amsmath,amsthm,amsfonts}
\usepackage[shortlabels]{enumitem}
\usepackage{multicol,multirow}
\usepackage{lipsum}
\usepackage{balance}
\usepackage{calc}
\usepackage[colorlinks=true,citecolor=blue,linkcolor=blue]{hyperref}
\usepackage{import}
\usepackage{xifthen}
\usepackage{pdfpages}
\usepackage{transparent}
\usepackage{tabularx}
\usepackage{listings}
\usepackage[dvipsnames]{xcolor}
\usepackage{bbm}

\newcommand{\incfig}[2][1.0]{
	\def\svgwidth{#1\columnwidth}
	\import{./figures/}{#2.pdf_tex}
}
\newcommand{\incimg}[2][1.0]{
	\includegraphics[width=#1\columnwidth]{./figures/#2}
}

\newcommand{\xor}{\oplus}
\newcommand{\floor}[1]{\left\lfloor #1 \right\rfloor}
\newcommand{\ceil}[1]{\left\lceil #1 \right\rceil}

\newlist{enumproof}{enumerate}{4}
\setlist[enumproof,1]{label=\arabic*., parsep=1em}
\setlist[enumproof,2]{label=\arabic{enumproofi}.\arabic*., parsep=1em}
\setlist[enumproof,3]{label=\arabic{enumproofi}.\arabic{enumproofii}.\arabic*., parsep=1em}
\setlist[enumproof,4]{label=\arabic{enumproofi}.\arabic{enumproofii}.\arabic{enumproofiii}.\arabic*., parsep=1em}

% ================= Styles listings =================
\definecolor{codegray}{rgb}{0.5,0.5,0.5}
\definecolor{codegreen}{rgb}{0,0.6,0}
\definecolor{mauve}{rgb}{0.58,0,0.82}
\definecolor{identblue}{rgb}{0.1,0.3,0.7}
\lstdefinestyle{base}{
	commentstyle=\color{codegray}\itshape,
	keywordstyle=\color{blue},
	stringstyle=\color{mauve},
	basicstyle={\ttfamily\footnotesize},
	breakatwhitespace=true,
	breaklines=true,
	captionpos=b,
	keepspaces=true,
	numbers=none,
	tabsize=2,
	showstringspaces=false,
}

% ---- Assembleur x86-64 (GAS / AT&T) ----
\lstdefinelanguage{gas}{
	sensitive=true,
	alsoletter={.\%_},
	morecomment=[l]{\#},
	morestring=[b]",
	% instructions
	morekeywords=[1]{mov,movb,movw,movl,movq,movabsq,movzbl,movzbq,movzwl,movzwq,
			movsbl,movsbq,movswl,movswq,movslq,cltq,cqto,lea,leaq,push,pushq,pop,popq,
			add,addl,addq,sub,subl,subq,imul,imull,imulq,mul,mulq,idiv,idivl,idivq,div,divq,
			inc,incl,incq,dec,decl,decq,neg,negl,negq,not,notl,notq,and,andl,andq,or,orl,orq,
			xor,xorl,xorq,sal,salq,sar,sarq,shl,shlq,shr,shrq,cmp,cmpl,cmpq,test,testl,testq,
			jmp,je,jne,jz,jnz,jl,jle,jg,jge,ja,jae,jb,jbe,js,jns,
			sete,setne,setz,setnz,setl,setle,setg,setge,seta,setae,setb,setbe,
			cmove,cmovne,cmovl,cmovle,cmovg,cmovge,call,ret,leave,enter,syscall},
	% directives
	morekeywords=[2]{.text,.data,.bss,.section,.rodata,.globl,.quad,.long,.word,
			.byte,.ascii,.asciz,.string,.file,.type,.size},
	% registres
	morekeywords=[3]{\%rax,\%rbx,\%rcx,\%rdx,\%rsi,\%rdi,\%rsp,\%rbp,\%rip,
			\%eax,\%ebx,\%ecx,\%edx,\%esi,\%edi,\%esp,\%ebp,
			\%ax,\%bx,\%cx,\%dx,\%al,\%bl,\%cl,\%dl,\%ah,\%bh,\%ch,\%dh,\%sil,\%dil,
			\%r8,\%r9,\%r10,\%r11,\%r12,\%r13,\%r14,\%r15,\%r8d,\%r9d,\%r8b,\%r9b},
}
\lstdefinestyle{asm}{
	style=base, language=gas,
	keywordstyle=[1]\color{blue},        % instructions
	keywordstyle=[2]\color{mauve},       % directives
	keywordstyle=[3]\color{codegreen},   % registres
}

% ---- C ----
\lstdefinestyle{c}{
	style=base, language=C,
	morekeywords=[2]{malloc,free,printf,scanf,sizeof,print_int,print_string},
	keywordstyle=[2]\color{identblue},   % libC / fonctions du projet
}

% ---- OCaml (+ ocamllex, Menhir, DSL du TP) ----
\lstdefinestyle{ocaml}{
	style=base, language=Caml,
	alsoletter={\%},
	morekeywords=[2]{rule,parse,eof,lexbuf,\%token,\%start,\%type},
	keywordstyle=[2]\color{identblue},   % ocamllex / Menhir
	morekeywords=[3]{zero,un,nouvelle_tension,nouveau_bus,relie,relie_bus,
			porte_nand,porte_delay,ins_cst_low,ins_cst_hi,ins_output,ins_jump_zero,
			ins_add,ins_inc,ins_load,ins_store},
	keywordstyle=[3]\color{codegreen},   % DSL circuits / CPU
}

% ---- Pseudo-code (descente récursive, algo LR) ----
\lstdefinelanguage{pseudo}{
	sensitive=true,
	morekeywords={def,if,else,return,while,for,in,loop,match,assert},
	morecomment=[l]{\#},
	morecomment=[s]{(*}{*)},
	morestring=[b]',
}
\lstdefinestyle{pseudo}{
	style=base, language=pseudo,
	morekeywords=[2]{peek,eat,push,pop,top,Shift,Reduce,Accept},
	keywordstyle=[2]\color{identblue},
}

% Style par défaut (\lstinline et listings sans option)
\lstset{style=base}

\usepackage{ifthen}
\usepackage[landscape]{geometry}
\usepackage[shortlabels]{enumitem}

\ifthenelse{\lengthtest { \paperwidth = 11in}}
    { \geometry{top=.5in,left=.5in,right=.5in,bottom=.5in} }
	{\ifthenelse{ \lengthtest{ \paperwidth = 297mm}}
		{\geometry{top=1cm,left=1cm,right=1cm,bottom=1cm} }
		{\geometry{top=1cm,left=1cm,right=1cm,bottom=1cm} }
	}

\pagestyle{empty}
\makeatletter
\renewcommand\thesection{\arabic{section}.}
\renewcommand{\section}{\@startsection{section}{1}{0mm}%
                                {-1ex plus -.5ex minus -.2ex}%
                                {0.05ex}%x
                                {\normalfont\normalsize\bfseries}}
\renewcommand{\subsection}{\@startsection{subsection}{2}{0mm}%
                                {-1ex plus -.5ex minus -.2ex}%
                                {0.05ex}%
                                {\normalfont\small\bfseries}}
\renewcommand{\subsubsection}{\@startsection{subsubsection}{3}{0mm}%
                                {-1ex plus -.5ex minus -.2ex}%
                                {0.05ex}%
                                {\normalfont\footnotesize\bfseries}}
\newcommand{\colbreak}{\vfill\null\columnbreak}
\makeatother
\setcounter{secnumdepth}{1}
\setlength{\parindent}{0pt}
\setlength{\parskip}{0.7em}

\setlist[itemize]{itemsep=0.6ex, topsep=2pt, partopsep=0pt, parsep=0pt, leftmargin=*, label=\textbullet}
\setlist[enumerate]{itemsep=0.6ex, topsep=-2pt, partopsep=0pt, parsep=0pt, label=\roman*.}

% Things Lie
\newcommand{\kb}{\mathfrak b}
\newcommand{\kg}{\mathfrak g}
\newcommand{\kh}{\mathfrak h}
\newcommand{\kn}{\mathfrak n}
\newcommand{\ku}{\mathfrak u}
\newcommand{\kz}{\mathfrak z}
\DeclareMathOperator{\Ext}{Ext} % Ext functor
\DeclareMathOperator{\Tor}{Tor} % Tor functor
\newcommand{\gl}{\opname{\mathfrak{gl}}} % frak gl group
\renewcommand{\sl}{\opname{\mathfrak{sl}}} % frak sl group chktex 6

% More script letters etc.
\newcommand{\SA}{\mathcal A}
\newcommand{\SB}{\mathcal B}
\newcommand{\SC}{\mathcal C}
\newcommand{\SF}{\mathcal F}
\newcommand{\SG}{\mathcal G}
\newcommand{\SH}{\mathcal H}
\newcommand{\OO}{\mathcal O}

\newcommand{\SCA}{\mathscr A}
\newcommand{\SCB}{\mathscr B}
\newcommand{\SCC}{\mathscr C}
\newcommand{\SCD}{\mathscr D}
\newcommand{\SCE}{\mathscr E}
\newcommand{\SCF}{\mathscr F}
\newcommand{\SCG}{\mathscr G}
\newcommand{\SCH}{\mathscr H}

% Mathfrak primes
\newcommand{\km}{\mathfrak m}
\newcommand{\kp}{\mathfrak p}
\newcommand{\kq}{\mathfrak q}

% number sets
\newcommand{\RR}[1][]{\ensuremath{\ifstrempty{#1}{\mathbb{R}}{\mathbb{R}^{#1}}}}
\newcommand{\NN}[1][]{\ensuremath{\ifstrempty{#1}{\mathbb{N}}{\mathbb{N}^{#1}}}}
\newcommand{\ZZ}[1][]{\ensuremath{\ifstrempty{#1}{\mathbb{Z}}{\mathbb{Z}^{#1}}}}
\newcommand{\QQ}[1][]{\ensuremath{\ifstrempty{#1}{\mathbb{Q}}{\mathbb{Q}^{#1}}}}
\newcommand{\CC}[1][]{\ensuremath{\ifstrempty{#1}{\mathbb{C}}{\mathbb{C}^{#1}}}}
\newcommand{\PP}[1][]{\ensuremath{\ifstrempty{#1}{\mathbb{P}}{\mathbb{P}^{#1}}}}
\newcommand{\HH}[1][]{\ensuremath{\ifstrempty{#1}{\mathbb{H}}{\mathbb{H}^{#1}}}}
\newcommand{\FF}[1][]{\ensuremath{\ifstrempty{#1}{\mathbb{F}}{\mathbb{F}^{#1}}}}
% expected value
\newcommand{\EE}{\ensuremath{\mathbb{E}}}
\newcommand{\charin}{\text{ char }}
\DeclareMathOperator{\sign}{sign}
\DeclareMathOperator{\Aut}{Aut}
\DeclareMathOperator{\Inn}{Inn}
\DeclareMathOperator{\Syl}{Syl}
\DeclareMathOperator{\Gal}{Gal}
\DeclareMathOperator{\GL}{GL} % General linear group
\DeclareMathOperator{\SL}{SL} % Special linear group

%---------------------------------------
% BlackBoard Math Fonts :-
%---------------------------------------

%Captital Letters
\newcommand{\bbA}{\mathbb{A}}	\newcommand{\bbB}{\mathbb{B}}
\newcommand{\bbC}{\mathbb{C}}	\newcommand{\bbD}{\mathbb{D}}
\newcommand{\bbE}{\mathbb{E}}	\newcommand{\bbF}{\mathbb{F}}
\newcommand{\bbG}{\mathbb{G}}	\newcommand{\bbH}{\mathbb{H}}
\newcommand{\bbI}{\mathbb{I}}	\newcommand{\bbJ}{\mathbb{J}}
\newcommand{\bbK}{\mathbb{K}}	\newcommand{\bbL}{\mathbb{L}}
\newcommand{\bbM}{\mathbb{M}}	\newcommand{\bbN}{\mathbb{N}}
\newcommand{\bbO}{\mathbb{O}}	\newcommand{\bbP}{\mathbb{P}}
\newcommand{\bbQ}{\mathbb{Q}}	\newcommand{\bbR}{\mathbb{R}}
\newcommand{\bbS}{\mathbb{S}}	\newcommand{\bbT}{\mathbb{T}}
\newcommand{\bbU}{\mathbb{U}}	\newcommand{\bbV}{\mathbb{V}}
\newcommand{\bbW}{\mathbb{W}}	\newcommand{\bbX}{\mathbb{X}}
\newcommand{\bbY}{\mathbb{Y}}	\newcommand{\bbZ}{\mathbb{Z}}

%---------------------------------------
% MathCal Fonts :-
%---------------------------------------

%Captital Letters
\newcommand{\mcA}{\mathcal{A}}	\newcommand{\mcB}{\mathcal{B}}
\newcommand{\mcC}{\mathcal{C}}	\newcommand{\mcD}{\mathcal{D}}
\newcommand{\mcE}{\mathcal{E}}	\newcommand{\mcF}{\mathcal{F}}
\newcommand{\mcG}{\mathcal{G}}	\newcommand{\mcH}{\mathcal{H}}
\newcommand{\mcI}{\mathcal{I}}	\newcommand{\mcJ}{\mathcal{J}}
\newcommand{\mcK}{\mathcal{K}}	\newcommand{\mcL}{\mathcal{L}}
\newcommand{\mcM}{\mathcal{M}}	\newcommand{\mcN}{\mathcal{N}}
\newcommand{\mcO}{\mathcal{O}}	\newcommand{\mcP}{\mathcal{P}}
\newcommand{\mcQ}{\mathcal{Q}}	\newcommand{\mcR}{\mathcal{R}}
\newcommand{\mcS}{\mathcal{S}}	\newcommand{\mcT}{\mathcal{T}}
\newcommand{\mcU}{\mathcal{U}}	\newcommand{\mcV}{\mathcal{V}}
\newcommand{\mcW}{\mathcal{W}}	\newcommand{\mcX}{\mathcal{X}}
\newcommand{\mcY}{\mathcal{Y}}	\newcommand{\mcZ}{\mathcal{Z}}


%---------------------------------------
% Bold Math Fonts :-
%---------------------------------------

%Captital Letters
\newcommand{\bmA}{\boldsymbol{A}}	\newcommand{\bmB}{\boldsymbol{B}}
\newcommand{\bmC}{\boldsymbol{C}}	\newcommand{\bmD}{\boldsymbol{D}}
\newcommand{\bmE}{\boldsymbol{E}}	\newcommand{\bmF}{\boldsymbol{F}}
\newcommand{\bmG}{\boldsymbol{G}}	\newcommand{\bmH}{\boldsymbol{H}}
\newcommand{\bmI}{\boldsymbol{I}}	\newcommand{\bmJ}{\boldsymbol{J}}
\newcommand{\bmK}{\boldsymbol{K}}	\newcommand{\bmL}{\boldsymbol{L}}
\newcommand{\bmM}{\boldsymbol{M}}	\newcommand{\bmN}{\boldsymbol{N}}
\newcommand{\bmO}{\boldsymbol{O}}	\newcommand{\bmP}{\boldsymbol{P}}
\newcommand{\bmQ}{\boldsymbol{Q}}	\newcommand{\bmR}{\boldsymbol{R}}
\newcommand{\bmS}{\boldsymbol{S}}	\newcommand{\bmT}{\boldsymbol{T}}
\newcommand{\bmU}{\boldsymbol{U}}	\newcommand{\bmV}{\boldsymbol{V}}
\newcommand{\bmW}{\boldsymbol{W}}	\newcommand{\bmX}{\boldsymbol{X}}
\newcommand{\bmY}{\boldsymbol{Y}}	\newcommand{\bmZ}{\boldsymbol{Z}}
%Small Letters
\newcommand{\bma}{\boldsymbol{a}}	\newcommand{\bmb}{\boldsymbol{b}}
\newcommand{\bmc}{\boldsymbol{c}}	\newcommand{\bmd}{\boldsymbol{d}}
\newcommand{\bme}{\boldsymbol{e}}	\newcommand{\bmf}{\boldsymbol{f}}
\newcommand{\bmg}{\boldsymbol{g}}	\newcommand{\bmh}{\boldsymbol{h}}
\newcommand{\bmi}{\boldsymbol{i}}	\newcommand{\bmj}{\boldsymbol{j}}
\newcommand{\bmk}{\boldsymbol{k}}	\newcommand{\bml}{\boldsymbol{l}}
\newcommand{\bmm}{\boldsymbol{m}}	\newcommand{\bmn}{\boldsymbol{n}}
\newcommand{\bmo}{\boldsymbol{o}}	\newcommand{\bmp}{\boldsymbol{p}}
\newcommand{\bmq}{\boldsymbol{q}}	\newcommand{\bmr}{\boldsymbol{r}}
\newcommand{\bms}{\boldsymbol{s}}	\newcommand{\bmt}{\boldsymbol{t}}
\newcommand{\bmu}{\boldsymbol{u}}	\newcommand{\bmv}{\boldsymbol{v}}
\newcommand{\bmw}{\boldsymbol{w}}	\newcommand{\bmx}{\boldsymbol{x}}
\newcommand{\bmy}{\boldsymbol{y}}	\newcommand{\bmz}{\boldsymbol{z}}

%---------------------------------------
% Scr Math Fonts :-
%---------------------------------------

\newcommand{\sA}{{\mathscr{A}}}   \newcommand{\sB}{{\mathscr{B}}}
\newcommand{\sC}{{\mathscr{C}}}   \newcommand{\sD}{{\mathscr{D}}}
\newcommand{\sE}{{\mathscr{E}}}   \newcommand{\sF}{{\mathscr{F}}}
\newcommand{\sG}{{\mathscr{G}}}   \newcommand{\sH}{{\mathscr{H}}}
\newcommand{\sI}{{\mathscr{I}}}   \newcommand{\sJ}{{\mathscr{J}}}
\newcommand{\sK}{{\mathscr{K}}}   \newcommand{\sL}{{\mathscr{L}}}
\newcommand{\sM}{{\mathscr{M}}}   \newcommand{\sN}{{\mathscr{N}}}
\newcommand{\sO}{{\mathscr{O}}}   \newcommand{\sP}{{\mathscr{P}}}
\newcommand{\sQ}{{\mathscr{Q}}}   \newcommand{\sR}{{\mathscr{R}}}
\newcommand{\sS}{{\mathscr{S}}}   \newcommand{\sT}{{\mathscr{T}}}
\newcommand{\sU}{{\mathscr{U}}}   \newcommand{\sV}{{\mathscr{V}}}
\newcommand{\sW}{{\mathscr{W}}}   \newcommand{\sX}{{\mathscr{X}}}
\newcommand{\sY}{{\mathscr{Y}}}   \newcommand{\sZ}{{\mathscr{Z}}}


%---------------------------------------
% Math Fraktur Font
%---------------------------------------

%Captital Letters
\newcommand{\mfA}{\mathfrak{A}}	\newcommand{\mfB}{\mathfrak{B}}
\newcommand{\mfC}{\mathfrak{C}}	\newcommand{\mfD}{\mathfrak{D}}
\newcommand{\mfE}{\mathfrak{E}}	\newcommand{\mfF}{\mathfrak{F}}
\newcommand{\mfG}{\mathfrak{G}}	\newcommand{\mfH}{\mathfrak{H}}
\newcommand{\mfI}{\mathfrak{I}}	\newcommand{\mfJ}{\mathfrak{J}}
\newcommand{\mfK}{\mathfrak{K}}	\newcommand{\mfL}{\mathfrak{L}}
\newcommand{\mfM}{\mathfrak{M}}	\newcommand{\mfN}{\mathfrak{N}}
\newcommand{\mfO}{\mathfrak{O}}	\newcommand{\mfP}{\mathfrak{P}}
\newcommand{\mfQ}{\mathfrak{Q}}	\newcommand{\mfR}{\mathfrak{R}}
\newcommand{\mfS}{\mathfrak{S}}	\newcommand{\mfT}{\mathfrak{T}}
\newcommand{\mfU}{\mathfrak{U}}	\newcommand{\mfV}{\mathfrak{V}}
\newcommand{\mfW}{\mathfrak{W}}	\newcommand{\mfX}{\mathfrak{X}}
\newcommand{\mfY}{\mathfrak{Y}}	\newcommand{\mfZ}{\mathfrak{Z}}
%Small Letters
\newcommand{\mfa}{\mathfrak{a}}	\newcommand{\mfb}{\mathfrak{b}}
\newcommand{\mfc}{\mathfrak{c}}	\newcommand{\mfd}{\mathfrak{d}}
\newcommand{\mfe}{\mathfrak{e}}	\newcommand{\mff}{\mathfrak{f}}
\newcommand{\mfg}{\mathfrak{g}}	\newcommand{\mfh}{\mathfrak{h}}
\newcommand{\mfi}{\mathfrak{i}}	\newcommand{\mfj}{\mathfrak{j}}
\newcommand{\mfk}{\mathfrak{k}}	\newcommand{\mfl}{\mathfrak{l}}
\newcommand{\mfm}{\mathfrak{m}}	\newcommand{\mfn}{\mathfrak{n}}
\newcommand{\mfo}{\mathfrak{o}}	\newcommand{\mfp}{\mathfrak{p}}
\newcommand{\mfq}{\mathfrak{q}}	\newcommand{\mfr}{\mathfrak{r}}
\newcommand{\mfs}{\mathfrak{s}}	\newcommand{\mft}{\mathfrak{t}}
\newcommand{\mfu}{\mathfrak{u}}	\newcommand{\mfv}{\mathfrak{v}}
\newcommand{\mfw}{\mathfrak{w}}	\newcommand{\mfx}{\mathfrak{x}}
\newcommand{\mfy}{\mathfrak{y}}	\newcommand{\mfz}{\mathfrak{z}}


\newcommand{\mytitle}{ECE\_3TC32\_TP Compilation}
\newcommand{\myauthor}{github/omgeta}
\newcommand{\mydate}{AU 26-27 S1}

\newcommand{\Null}{\operatorname{Null}}
\newcommand{\Prem}{\operatorname{Premier}}
\newcommand{\Suiv}{\operatorname{Suivant}}
\newcommand{\Clo}{\textrm{Clôture}}

\begin{document}
\raggedright
\raggedcolumns
\footnotesize
\begin{multicols*}{3}
	\setlength{\premulticols}{1pt}
	\setlength{\postmulticols}{1pt}
	\setlength{\multicolsep}{1pt}
	\setlength{\columnsep}{2pt}

	{\normalsize{\textbf{\mytitle}}}\\
	{\footnotesize{\mydate\hspace{2pt}\textemdash\hspace{2pt}\myauthor}}
	%%%%%%%%%%%%%%%%%%%%%%%%%%%%%%%%%%%%%%%%%%%%%%%%%%%%%%
	%                      Begin                         %
	%%%%%%%%%%%%%%%%%%%%%%%%%%%%%%%%%%%%%%%%%%%%%%%%%%%%%%
	\section{Introduction}

	Compilateur est un programme qui traduit un "programme" d'un langage à un autre en préservant la sémantique (correct, rapide, optimisant).

	Grandes étapes:
	\begin{enumerate}[\roman*.]
		\item Analyse lexicale/Lexing:\\texte $\rightarrow$ tokens via expressions régulières
		\item Analyse syntaxique/Parsing:\\tokens $\rightarrow$ AST via une grammaire
		\item Analyse sémantique: variables/fonctions déclarées, types corrects, annotations de l'AST
		\item Optimisation
		\item Géneration de code:\\assembleur $\rightarrow$ object $\rightarrow$ binaire
	\end{enumerate}

	Representation de base-R : $\displaystyle m = \sum^{n-1}_{i=0} b_i \cdot R^i$

	Complément à deux:
	\begin{enumerate}[\roman*.]
		\item $b_{N-1}\ldots b_0 = -2^{N-1}b_{N-1}+\sum_{i<N-1}b_i2^i$
		\item Opposé de $x$: $-x = 2^N - x$
		\item Intervalle: $-2^{N-1}$ à $2^{N-1}-1$
		\item Négation: $-x = \sim x+1$
		\item Débordement: tronquer (calculs dans $\ZZ/2^N\ZZ$)
		\item \lstinline|+ - *| identiques signé/non signé
	\end{enumerate}

	Boutisme:
	\begin{enumerate}[\roman*.]
		\item Grand-boutien: écriture «naturelle»
		\item Petit-boutien: octet de poids faible à la basse; $m=o_1o_2 \Rightarrow m=o_1+256\,o_2$
	\end{enumerate}
	\vspace{1em}
	\begin{tabular}{|l|l|l|l|}
		\hline
		\textbf{Unité} & \textbf{Taille} & \textbf{C}                             & \textbf{courant} \\
		\hline
		octet/byte (b) & 8 bits          & \texttt{char}                          & 8                \\
		mot/word (w)   & 2 o             & \texttt{short}                         & 16               \\
		double (l)     & 4 o             & \texttt{int}                           & 32               \\
		quad (q)       & 8 o             & \texttt{long}, \texttt{long long}, ptr & 64               \\
		\hline
	\end{tabular}

	\section{x86-64}
	\lstset{style=asm}
	Modèle:
	\begin{enumerate}[\roman*.]
		\item Registres: rapides, peu nombreux
		\item RAM: grand tableau
		\item \lstinline|%rip|: adresse de l'instruction courante
		\item Instruction lit/écrit registres ou RAM, ou \lstinline|syscall|
	\end{enumerate}

	Syntaxe AT\&T: \lstinline|op source, destination|
	\begin{enumerate}[\roman*.]
		\item Registres préfixes par \lstinline|%| (e.g. \lstinline|$r9|)
		\item Constants préfixes par \lstinline|$| (e.g. \lstinline|$0x2A|) ;\\que les 32 bits signés
		\item Suffixe de taille: \lstinline|b|=1, \lstinline|w|=2, \lstinline|l|=4, \lstinline|q|=8 octets; obligatoire si ambigu (e.g. \lstinline|movq $42, (%rsp)|)
	\end{enumerate}

	\subsection{Registres Généraux}
	\begin{tabular}{|l|l|l|l|l|}
		\hline
		\textbf{64}               & \textbf{32}  & \textbf{16}  & \textbf{8}                & \textbf{Rôle historique } \\
		\hline
		\texttt{rax}              & \texttt{eax} & \texttt{ax}  & \texttt{al} (\texttt{ah}) & Accumulator / retour      \\
		\texttt{rbx}              & \texttt{ebx} & \texttt{bx}  & \texttt{bl} (\texttt{bh}) & Base / callee-saved       \\
		\texttt{rcx}              & \texttt{ecx} & \texttt{cx}  & \texttt{cl} (\texttt{ch}) & Counter / arg 4           \\
		\texttt{rdx}              & \texttt{edx} & \texttt{dx}  & \texttt{dl} (\texttt{dh}) & Data / arg 3              \\
		\texttt{rsi}              & \texttt{esi} & \texttt{si}  & \texttt{sil}              & Source idx / arg 2        \\
		\texttt{rdi}              & \texttt{edi} & \texttt{di}  & \texttt{dil}              & Dest. idx / arg 1         \\
		\texttt{rsp}              & \texttt{esp} & \texttt{sp}  & \texttt{spl}              & Stack Pointer             \\
		\texttt{rbp}              & \texttt{ebp} & \texttt{bp}  & \texttt{bpl}              & Base Pointer (frame)      \\
		\texttt{r8}--\texttt{r15} & \texttt{r8d} & \texttt{r8w} & \texttt{r8b}              & libres (r8,r9: args 5,6)  \\
		\hline
	\end{tabular}

	Addressage memoire:
	\begin{enumerate}[\roman*.]
		\item Schema: \lstinline|offset(base, index, scale)| \\$=$ offset $+$ base $+$ index$\times$scale
		\item offset: constante 32 bits signée; scale $\in\{1,2,4,8\}$
		\item \lstinline|mov x(%rip), %rsi|: contenu au label \lstinline|x|
		\item \lstinline|lea x(%rip), %rsi|: adresse de \lstinline|x|
	\end{enumerate}
	Mémoire virtuelle: tableau de $2^{64}$ octets (48 bits utilisés), traduit par pages par l'OS, avec droits R/W/X.

	\colbreak
	\subsection{Sections et données}
	\begin{lstlisting}
.text      # code (executable)
.data      # donnees initialisees
.bss       # donnees init 0
.rodata    # lecture seule
.globl main  # exporte le symbole l'ELF
w:  .quad 0x1234567890  # 8o
x:  .long 12        # 4o 
y:  .word 12        # 2o
z:  .byte 12      # 1o
s:  .ascii "Hello " # 6o
t:  .asciz "World"  # 6o, avec \0
\end{lstlisting}

	\subsection{Instructions}
	Copie:
	\begin{enumerate}[\roman*.]
		\item \lstinline|mov[bwlq] s, d|: \lstinline|d = s|
		\item \lstinline|movz[bw][wlq] s, d|: extension par zéros
		\item \lstinline!movs[bw|bl|bq|wl|wq|lq] s, d!: extension signée
		\item \lstinline|cqto|: étend le signe de \lstinline|%rax| dans \lstinline|%rdx:%rax| (avant \lstinline|idivq|)
		\item \lstinline|push s| / \lstinline|pop d|: modifier pile
	\end{enumerate}

	Arithmétique:
	\begin{enumerate}[\roman*.]
		\item \lstinline|add s, d| : \lstinline|d += s| / \lstinline|sub s, d|: \lstinline|d -= s|
		\item \lstinline|inc d| / \lstinline|dec d| / \lstinline|neg d|
		\item \lstinline|imul s, d|: \lstinline|d *= s| (mod $2^{64}$)
		\item \lstinline|imul s| : \lstinline|%rdx:%rax = %rax * s|; \lstinline|mul s| non signé
		\item \lstinline|idiv s|: divise \lstinline|%rdx:%rax| par \lstinline|s|; \lstinline|div s| non signé; \\quotient $\rightarrow$ \lstinline|%rax|, reste $\rightarrow$ \lstinline|%rdx|.
		\item \lstinline|lea s, d|: \lstinline|d = &s|
	\end{enumerate}

	Bit-à-bit:
	\begin{enumerate}[\roman*.]
		\item \lstinline|not d| / \lstinline|and s, d| / \lstinline|or s, d| / \lstinline|xor s, d|
		\item \lstinline|sal k, d|: \lstinline|d <<= k| / \lstinline|sar k, d|: \lstinline|d >>= k| signé
		\item \lstinline|shr k, d|: \lstinline|d >>= k| logique
	\end{enumerate}

	Comparaisons (flags uniquement):
	\begin{enumerate}[\roman*.]
		\item \lstinline|cmp s, d|: \lstinline|d - s|
		\item \lstinline|test s, d|: \lstinline|d & s|
	\end{enumerate}

	\colbreak

	Flags:
	\begin{enumerate}[\roman*.]
		\item \textbf{CF} retenue / \textbf{ZF} nul / \textbf{SF} négatif /\\ \textbf{OF} débordement signé / \textbf{PF} parité
		\item Modifiés par arithmétique, bit-à-bit, \lstinline|cmp|, \lstinline|test|
		\item Pas modifiés par \lstinline|mov|, \lstinline|lea|, \lstinline|push|/\lstinline|pop|, \lstinline|jCC|, \lstinline|setCC|
	\end{enumerate}

	Sauts:
	\begin{enumerate}[\roman*.]
		\item \lstinline|jmp label| / \lstinline|jmp *%rax|: inconditionnel
		\item \lstinline|jCC label| / \lstinline|cmovCC s, d|: si CC
		\item \lstinline|setCC d|: \lstinline|d = CC ? 1 : 0| (\lstinline|d| 8 bits, ex. \lstinline|%al|)
	\end{enumerate}

	\subsection{Conditions}
	\begin{tabular}{|l|l|l|}
		\hline
		\textbf{CC}               & \textbf{Sens}          & \textbf{Flags}                      \\
		\hline
		\texttt{e} / \texttt{z}   & $d = s$ / nul          & ZF                                  \\
		\texttt{ne} / \texttt{nz} & $d \neq s$             & $\lnot$ZF                           \\
		\texttt{s} / \texttt{ns}  & négatif / positif      & SF                                  \\
		\texttt{p} / \texttt{np}  & parité paire / impaire & PF                                  \\
		\hline
		\texttt{g} / \texttt{nle} & $d > s$ signé          & $\lnot$(SF$\oplus$OF)$\land\lnot$ZF \\
		\texttt{ge} / \texttt{nl} & $d \geq s$ signé       & $\lnot$(SF$\oplus$OF)               \\
		\texttt{l} / \texttt{nge} & $d < s$ signé          & SF$\oplus$OF                        \\
		\texttt{le} / \texttt{ng} & $d \leq s$ signé       & (SF$\oplus$OF)$\lor$ZF              \\
		\hline
		\texttt{a} / \texttt{nbe} & $d > s$ non signé      & $\lnot$CF$\land\lnot$ZF             \\
		\texttt{ae} / \texttt{nb} & $d \geq s$ non signé   & $\lnot$CF                           \\
		\texttt{b} / \texttt{nae} & $d < s$ non signé      & CF                                  \\
		\texttt{be} / \texttt{na} & $d \leq s$ non signé   & CF$\lor$ZF                          \\
		\hline
	\end{tabular}


	\subsection{\lstinline|syscall| (Linux)}

	\begin{tabular}{|l|l|l|l|l|}
		\hline
		\texttt{\%rax} & \textbf{Nom} & \texttt{\%rdi} & \texttt{\%rsi} & \texttt{\%rdx} \\
		\hline
		0              & read         & fd             & \texttt{char*} & count          \\
		1              & write        & fd             & \texttt{char*} & count          \\
		12             & brk          & brk            &                &                \\
		60             & exit         & code           &                &                \\
		\hline
	\end{tabular}

	avec file descriptors: 0 = stdin, 1 = stdout, 2 = stderr.

	\colbreak

	\section{Mémoire et fonctions}
	Memoire d'un programme:
	\begin{enumerate}[\roman*.]
		\item Statique: \lstinline|.text/.data/.bss/.rodata|,\\chargée depuis l'ELF
		\item {Tas} (heap): à la demande (\lstinline|malloc|), croît $\uparrow$
		\item {Pile} (stack): fournie par l'OS, en haut, croît $\downarrow$
	\end{enumerate}

	Pile:
	\begin{enumerate}[\roman*.]
		\item \lstinline|push s|: \lstinline|sub $8,%rsp| + \lstinline|mov s,(%rsp)|
		\item \lstinline|pop d|: \lstinline|mov (%rsp),d| + \lstinline|add $8,%rsp|
		\item \lstinline|call f|: \lstinline|push| adresse retour + \lstinline|jmp f|
		\item \lstinline|ret|: \lstinline|pop %rip|
		\item \lstinline|enter $N|:\\\lstinline|push %rbp| + \lstinline|mov %rsp,%rbp| + \lstinline|sub $N,%rsp|
		\item \lstinline|leave|: \lstinline|mov %rbp,%rsp| + \lstinline|pop %rbp|
	\end{enumerate}
	\begin{tabular}{|l|l|}
		\hline
		\textbf{Pile (frame)} & \textbf{Adresse}                        \\
		\hline
		args 7+ de l'appelant & \texttt{16(\%rbp)}, \texttt{24(\%rbp)}  \\
		adresse de retour     & \texttt{8(\%rbp)}                       \\
		ancien \texttt{\%rbp} & $\leftarrow$ \texttt{\%rbp}             \\
		variables locales     & \texttt{-8(\%rbp)}, \texttt{-16(\%rbp)} \\
		temporaires           & $\leftarrow$ \texttt{\%rsp} (sommet)    \\
		\hline
	\end{tabular}

	\subsection{Convention d'appel (System V)}
	\begin{enumerate}[\roman*.]
		\item Arguments: \lstinline|%rdi, %rsi, %rdx, %rcx, %r8, %r9|, puis pile
		\item Retour: \lstinline|%rax|
		\item Callee-saved: \lstinline|%rbx, %rbp, %r12|--\lstinline|%r15|
		\item Caller-saved: \lstinline|%rax, %rcx, %rdx, %rsi, %rdi, %r8|--\lstinline|%r11|
		\item Variadique (\lstinline|printf|, \lstinline|scanf|): \lstinline|%al| = 0
		\item Alignement: \lstinline|%rsp| multiple de 16 au \lstinline|call|\\($\equiv 8 \bmod 16$ à l'entrée, aligné après \lstinline|push %rbp|)
	\end{enumerate}

	\colbreak
	Squelette:
	\noindent\begin{minipage}{\linewidth}
		\begin{lstlisting}
f:  push %rbp          # frame appelant
    mov  %rsp, %rbp    # nouvelle frame
    sub  $16, %rsp     # locales (x16)
    mov  %rdi, -8(%rbp)  # arg 1
    ...                # result -> %rax
    leave
    ret
\end{lstlisting}
	\end{minipage}

	\subsection{libC}
	\begin{lstlisting}
    .data
fmt_d:  .string "%d\n"
fmt_in: .string "%d"
    .text
    .globl main
main:
    push %rbp
    mov  %rsp, %rbp
    sub  $16, %rsp

    # scanf("%d", &n)
    lea  -8(%rbp), %rsi 
    lea  fmt_in(%rip), %rdi
    xor  %eax, %eax
    call scanf
    movslq -8(%rbp), %rsi   # int 4o>8o

    # printf("%d\n", n)
    lea  fmt_d(%rip), %rdi  
    xor  %eax, %eax
    call printf
    xor  %eax, %eax         # return 0
    leave
    ret
\end{lstlisting}
	\colbreak

	\section{Rappels de C}
	C est bas niveau, impératif, à passage par valeur et statiquement typé.

	Types:
	\begin{enumerate}[\roman*.]
		\item Fixent la taille (\lstinline|sizeof(T)|) et le comportement
		\item Modificateurs: \lstinline|signed|, \lstinline|unsigned|, \lstinline|const|
		\item Constantes: \lstinline|42|, \lstinline|0xA2|, \lstinline|'H'|, \lstinline|42LL|
	\end{enumerate}

	Stockage:
	\begin{enumerate}[\roman*.]
		\item Vie courte: registre
		\item Globale:\\\lstinline|.data| (init.), \lstinline|.bss| (non init.), \lstinline|.rodata| (\lstinline|const|)
		\item Locale: pile (une copie par appel)
	\end{enumerate}

	Valeur gauche: a une adresse, peut être à gauche de \lstinline|=|
	\begin{enumerate}[\roman*.]
		\item E.g.: \lstinline|v|, \lstinline|*v|, \lstinline|*(v+2)|, \lstinline|*(f(v))|
		\item Non: \lstinline|42|, \lstinline|v+42|, \lstinline|f(v)|
	\end{enumerate}

	Pointeurs:
	\begin{enumerate}[\roman*.]
		\item \lstinline|T* p|: pointeur vers un \lstinline|T|
		\item \lstinline|&a|: adresse de la l-value \lstinline|a|
		\item \lstinline|*p|: objet pointé
	\end{enumerate}

	Arithmétique de pointeurs (\lstinline|T* p, p1, p2|):
	\begin{enumerate}[\roman*.]
		\item \lstinline|p + i| $= p + i \times$ \lstinline|sizeof(T)|
		\item \lstinline|p1 - p2| $= (p_1 - p_2)\,/$ \lstinline|sizeof(T)|
	\end{enumerate}

	Tableaux:
	\begin{enumerate}[\roman*.]
		\item \lstinline|t[i]| $=$ \lstinline|*(t+i)|; \lstinline|&t[0] == t|
		\item Statique: \lstinline|int t[10];|
		\item Dynamique: \lstinline|int* t = malloc(10*sizeof(int));|
		\item Linéarisé: \lstinline|int t[99][88][77]| \\$\Rightarrow$ \lstinline|t[i][j][k]| $=$ \lstinline|*(t + i*88*77 + j*77 + k)|
	\end{enumerate}

	Chaînes: tableau de \lstinline|char| terminé par \lstinline|'\0'|; \lstinline|"abc"| est un \lstinline|const char*| en \lstinline|.rodata|.

	libC:
	\begin{enumerate}[\roman*.]
		\item \lstinline|malloc(n)| / \lstinline|free(p)|
		\item \lstinline|printf(fmt, ...)| / \lstinline|scanf(fmt, ...)|
	\end{enumerate}
	\vspace{-1em}
	\colbreak
  \section{Grammaires}
	Grammaire formelle $G$:
	\begin{enumerate}[\roman*.]
		\item Terminaux $\Sigma$ (minuscules): lettres des mots
		\item Non-terminaux $N$ (majuscules): variables
		\item Symbole de départ $S$
		\item Règles $(\Sigma\cup N)^+\rightarrow(\Sigma\cup N)^*$
	\end{enumerate}
	Langage $L(G)$: mots de $\Sigma^*$ obtenus depuis $S$ en appliquant les règles.
 
	Notations:
	\begin{enumerate}[\roman*.]
		\item $X\rightarrow\alpha$: règle; $\alpha$ est une production de $X$
		\item $\alpha\Rightarrow\beta$: une étape (un NT remplacé)
		\item $\alpha\Rightarrow^*\beta$: plusieurs étapes (dérivation)
		\item $\alpha\Rightarrow_L\beta$: toujours le NT le plus à gauche
	\end{enumerate}
 
	Décider $w\in L(G)$:
	\begin{enumerate}[\roman*.]
		\item Grammaire quelconque: \textbf{indécidable}
		\item CFG (règles $N\rightarrow(\Sigma\cup N)^*$): $O(n^3|G|)$, trop lent
		\item LL($k$), LR($k$): \textbf{linéaire} $\Rightarrow$ utilisées en pratique
	\end{enumerate}
 
	Inclusions:
	\begin{enumerate}[\roman*.]
		\item LL(0) $\subset$ LL(1) $\subset$ LL($k$) $\subset$ LR($k$) $\subset$ CFG non ambiguës
		\item LR(0) $\subset$ SLR(1) $\subset$ LALR(1) $\subset$ LR(1)
	\end{enumerate}
 
	\subsection{Arbre de dérivation}
	\begin{enumerate}[\roman*.]
		\item Racine: $S$\\Nœuds internes: NT\\Feuilles: terminaux
		\item Enfants de $X$: $s_1\ldots s_k$ pour une règle $X\rightarrow s_1\ldots s_k$
		\item Mot de l'arbre: feuilles lues de gauche à droite
		\item $w\in L(G)\iff w$ est le mot d'un arbre
	\end{enumerate}
	\textbf{Ambiguë}: un mot a plusieurs arbres (le sens du programme n'est plus unique).
 
	\subsection{Lever l'ambiguïté}
	Dangling else: rattacher \lstinline|else| au \lstinline|if| le plus imbriqué ($IE$ = \lstinline|if| avec son \lstinline|else|):
	\begin{gather*}
		S \rightarrow \text{if } S \mid IE\\
		IE \rightarrow \text{printf} \mid \text{if } IE \text{ else } S
	\end{gather*}
	Précédence: un NT par niveau, du plus faible au plus fort:
	\begin{gather*}
		S_0 \rightarrow S_0 + S_1 \mid S_1\\
		S_1 \rightarrow S_1 * S_2 \mid S_2\\
		S_2 \rightarrow S_3\,\hat{}\,S_2 \mid S_3\\
		S_3 \rightarrow \text{int} \mid (S_0)
	\end{gather*}
	Le sens de la récursion fixe l'associativité:
	\begin{enumerate}[\roman*.]
		\item Récursion gauche ($S_0\rightarrow S_0+S_1$): gauche,\\$1-2-3=(1-2)-3$
		\item Récursion droite ($S_2\rightarrow S_3\,\hat{}\,S_2$): droite,\\$2\,\hat{}\,3\,\hat{}\,4=2^{(3^4)}$
	\end{enumerate}
 
	\subsection{Algorithme général (CFG)}
	\begin{enumerate}[\roman*.]
		\item Binariser: $X\rightarrow s_1\ldots s_k$ devient\\$X\rightarrow s_1F$, $F\rightarrow s_2\ldots s_k$ ($F$ frais)
		\item $(i,j,\alpha)\in match$ ssi $\alpha\Rightarrow^*w_i\ldots w_{j-1}$
		\item $w\in L(G)\iff(0,n,S)\in match$
	\end{enumerate}
	$match$: plus petit ensemble contenant:
	\begin{enumerate}[\roman*.]
		\item $(i,i,\alpha)$ si $\alpha\rightarrow\epsilon$; $(i,i+1,w_i)$
		\item $(i,j,\alpha)$ si $\alpha\rightarrow\beta$ et $(i,j,\beta)$
		\item $(i,j,\alpha)$ si $\alpha\rightarrow\beta\gamma$, $(i,k,\beta)$ et $(k,j,\gamma)$
	\end{enumerate}
 
  \colbreak
	\section{LL(1)}
	Analyse descendante: lecture \textbf{L}eft-to-right, dérivation \textbf{L}eftmost, $k$ symboles d'avance.\\
	LL(1): \textbf{le prochain symbole suffit à choisir la règle}.
 
	\subsection{Null, Premier, Suivant}
	Calculés par point fixe (répéter jusqu'à stabilité).
 
	$\Null$: NT $X$ avec $X\Rightarrow^*\epsilon$
	\begin{enumerate}[\roman*.]
		\item $X\rightarrow\epsilon$ $\Rightarrow$ $X\in\Null$
		\item $X\rightarrow Y_1\ldots Y_k$, tous $Y_i\in\Null$ $\Rightarrow$ $X\in\Null$
		\item Un terminal n'est jamais nul
	\end{enumerate}
 
	$\Prem(\alpha)$: terminaux qui commencent un mot dérivé de $\alpha$
	\begin{enumerate}[\roman*.]
		\item $\Prem(t)=\{t\}$ pour $t$ terminal
		\item $\Prem(X)=\bigcup_{X\rightarrow\alpha}\Prem(\alpha)$
		\item $\Prem(\alpha_1\alpha_2\ldots)=\Prem(\alpha_1)$,\\$\cup\ \Prem(\alpha_2\ldots)$ si $\alpha_1\in\Null$
	\end{enumerate}
 
	$\Suiv(X)$: terminaux qui suivent $X$ ($S\Rightarrow^*\alpha Xt\beta$)
	\begin{enumerate}[\roman*.]
		\item $\$\in\Suiv(S)$ (fin de l'entrée)
		\item Pour chaque $X\rightarrow\alpha Y\beta$: $\Prem(\beta)\subseteq\Suiv(Y)$
		\item Si de plus $\Null(\beta)$: $\Suiv(X)\subseteq\Suiv(Y)$
	\end{enumerate}
 
	\subsection{Tester LL(1)}
	\begin{enumerate}[\roman*.]
		\item Calculer $\Null$, puis $\Prem$, puis $\Suiv$
		\item Pour chaque règle $X\rightarrow\alpha$, la mettre dans $M[X,t]$:
		\begin{itemize}
			\item pour tout $t\in\Prem(\alpha)$
			\item pour tout $t\in\Suiv(X)$ si $\Null(\alpha)$
		\end{itemize}
		\item LL(1) $\iff$ au plus une règle par case
	\end{enumerate}
	Conflits pour $X\rightarrow\alpha\mid\beta$:
	\begin{enumerate}[\roman*.]
		\item $\Prem(\alpha)\cap\Prem(\beta)\neq\emptyset$
		\item $\alpha$ et $\beta$ dérivent toutes deux $\epsilon$
		\item $\alpha\Rightarrow^*\epsilon$ et $\Prem(\beta)\cap\Suiv(X)\neq\emptyset$
	\end{enumerate}
 
  \colbreak
	\subsection{Rendre LL(1)}
	Problèmes:
	\begin{enumerate}[\roman*.]
		\item Récursion gauche: $S\rightarrow S+T$ (boucle infinie)
		\item Préfixe commun: $S\rightarrow P\mid P+S$ (même $\Prem$)
	\end{enumerate}
	Solution: factoriser avec un NT de suite $F_i$.\\
	$S_0\rightarrow S_1+S_0\mid S_1$, $S_1\rightarrow S_2*S_1\mid S_2$, $S_2\rightarrow$ int devient:
	\begin{gather*}
		S_0 \rightarrow S_1F_0 \qquad F_0 \rightarrow +S_0 \mid \epsilon\\
		S_1 \rightarrow S_2F_1 \qquad F_1 \rightarrow *S_1 \mid \epsilon\\
		S_2 \rightarrow \text{int}
	\end{gather*}
	\begin{enumerate}[\roman*.]
		\item $\Null=\{F_0,F_1\}$
		\item $\Prem(S_i)=\{\text{int}\}$, $\Prem(F_0)=\{+\}$, $\Prem(F_1)=\{*\}$
		\item $\Suiv(S_0)=\Suiv(F_0)=\{\$\}$
		\item $\Suiv(S_1)=\Suiv(F_1)=\{+,\$\}$
		\item $\Suiv(S_2)=\{*,+,\$\}$
	\end{enumerate}
 
	\subsection{Descente récursive}
	\begin{enumerate}[\roman*.]
		\item Une fonction par NT, qui renvoie l'arbre
		\item \lstinline|peek()|: regarde le prochain symbole
		\item \lstinline|eat()|: le consomme
		\item Règle choisie selon \lstinline|peek()| (table LL(1))
	\end{enumerate}
	\noindent\begin{minipage}{\linewidth}
		\begin{lstlisting}[style=pseudo]
def s0():          # S0 -> S1 + S0 | S1
    a = s1()
    if peek() != '+': return a
    eat(); return (a, '+', s0())
def s1():          # S1 -> S2 * S1 | S2
    a = s2()
    if peek() != '*': return a
    eat(); return (a, '*', s1())
def s2():          # S2 -> int
    a = peek(); assert a == 'int'
    eat(); return a
\end{lstlisting}
	\end{minipage}
	Attention: associativité droite ($1-2-3=1-(2-3)$). Pour gauche, une boucle:\\\lstinline|while peek()=='-': eat(); a=(a,'-',s1())|
	\colbreak

  \section{Lexeur}
	Rôle: découper le texte en tokens, décrits par des expressions régulières.\\
	e.g. \lstinline|let x = 42| $\rightarrow$ LET VAR(x) EQ INT(42)
 
	Langages réguliers: CFG avec règles $N\rightarrow\Sigma N\mid\epsilon$\\(= regex = automates).
 
	\subsection{Expressions régulières}
	\begin{tabular}{|l|l|l|}
		\hline
		$r$          & \textbf{Sens} & $L(r)$                                 \\
		\hline
		$\emptyset$  & vide          & $\emptyset$                            \\
		$\epsilon$   & mot vide      & $\{\epsilon\}$                         \\
		$c$          & caractère     & $\{c\}$                                \\
		$[c_1-c_2]$  & classe        & $\{c\mid c_1\leq c\leq c_2\}$          \\
		$r_a|r_b$    & union         & $L(r_a)\cup L(r_b)$                    \\
		$r_ar_b$     & concaténation & $\{w_aw_b\}$                           \\
		$r^*$        & $\geq0$ fois  & $\{w_1\ldots w_k\mid k\geq0\}$         \\
		$r^+$        & $\geq1$ fois  & $\{w_1\ldots w_k\mid k\geq1\}$         \\
		\hline
	\end{tabular}
 
	\subsection{Algorithme du lexeur}
	Entrée: liste ordonnée (regex $e_i$, action $a_i$). Sur $w$:
	\begin{enumerate}[\roman*.]
		\item $w$ vide: fin
		\item Chercher le \textbf{plus long préfixe} reconnu par un $e_i$ (aucun: erreur)
		\item Égalité: \textbf{plus petit $i$} (\lstinline|let| avant \lstinline|[a-z]+|)
		\item Exécuter $a_i$ sur le préfixe, recommencer sur le reste
	\end{enumerate}
 
	Lexeur séparé: décompose le problème, grammaires plus simples.
 
  \colbreak
	\subsection{Reconnaître une regex}
	$match(R,i)=\{j\mid w_i\ldots w_{j-1}\in L(R)\}$:
	\begin{enumerate}[\roman*.]
		\item $match(\epsilon,i)=\{i\}$
		\item $match(c,i)=\{i+1\}$ si $w_i=c$
		\item $match([c_1,c_2],i)=\{i+1\}$ si $c_1\leq w_i\leq c_2$
		\item $match(R_a|R_b,i)=match(R_a,i)\cup match(R_b,i)$
		\item $match(R_aR_b,i)=\bigcup_{j\in match(R_a,i)}match(R_b,j)$
		\item $match(R^*,i)$: point fixe $X=\{i\}\cup\bigcup_{j\in X}match(R,j)$
	\end{enumerate}
 
	Automate fini déterministe:
	\begin{enumerate}[\roman*.]
		\item États $Q$, finaux $F\subseteq Q$, initial $i$
		\item Transition $\delta:Q\times\Sigma\rightarrow Q$
		\item $w$ accepté si $q_n\in F$, avec $q_0=i$, $q_{k+1}=\delta(q_k,w_{k+1})$
	\end{enumerate}
	En pratique: écrire des regex (lisibles), les compiler en automates (efficaces).
 
	\subsection{ocamllex}
	\noindent\begin{minipage}{\linewidth}
		\begin{lstlisting}[style=ocaml]
rule token = parse
| ['0'-'9']+ as s
          { `INT (int_of_string s) }
| "let"            { `LET }
| '+'              { `PLUS }
| [' ' '\t' '\n']+ { token lexbuf }
| ['a'-'z']['a'-'z' '0'-'9']* as id
                   { `VAR id }
| eof              { `EOF }
\end{lstlisting}
	\end{minipage}
 
  \colbreak
	\section{LR}
	Analyse ascendante: lecture \textbf{L}eft-to-right, dérivation \textbf{R}ightmost construite à l'envers, $k$ symboles d'avance. Plus puissant que LL (récursion gauche OK).
 
	\subsection{Shift-reduce}
	Une pile de symboles + l'entrée. À chaque étape:
	\begin{enumerate}[\roman*.]
		\item Shift: empiler le prochain symbole
		\item Reduce $A\rightarrow\alpha$: dépiler $\alpha$, empiler $A$
		\item Accepter / refuser
	\end{enumerate}
	L'action est lue dans une \textbf{table LR}: (état de la pile, prochain symbole) $\rightarrow$ action.
 
	\subsection{Analyser avec la table}
	\begin{enumerate}[\roman*.]
		\item Pile d'états (débute à $0$) + pile de symboles
		\item \textsc{Action}[état, symbole]: shift $s$, reduce, accept ou erreur
		\item Après reduce $X\rightarrow\alpha$: dépiler $|\alpha|$, empiler \textsc{Goto}[état, $X$]
		\item Arbre: chaque reduce crée le nœud $X$ avec les $|\alpha|$ fils dépilés
	\end{enumerate}
	\noindent\begin{minipage}{\linewidth}
		\begin{lstlisting}[style=pseudo]
states <- [0] ; symbols <- []
loop:
  match action(states.top, input.peek):
  | Shift(s) ->
      symbols.push(input.pop())
      states.push(s)
  | Reduce(X->a) ->
      pop |a| de states et symbols
      symbols.push(X)
      states.push(goto(states.top, X))
  | Accept -> return OK
  | vide   -> erreur de syntaxe
\end{lstlisting}
	\end{minipage}
 
  \colbreak
	\subsection{Construire la table LR(1)}
	Item LR(1) $[X\rightarrow\alpha\bullet\beta,\ c]$: $\alpha$ déjà reconnu, $\beta$ cherché, le tout suivi de $c$ (LR(0): sans $c$).
 
	Clôture: pour $[X\rightarrow\alpha\bullet Y\beta,c]\in I$, ajouter $[Y\rightarrow\bullet\gamma,d]$ pour toute règle $Y\rightarrow\gamma$ et tout $d\in\Prem(\beta c)$.
 
	Transition: $\delta(I,X)=\Clo$ des $[Y\rightarrow\alpha X\bullet\gamma,c]$ pour $[Y\rightarrow\alpha\bullet X\gamma,c]\in I$.
 
	Méthode:
	\begin{enumerate}[\roman*.]
		\item Augmenter: $G\rightarrow S\$$; $I_0=\Clo(\{[G\rightarrow\bullet S\$,\$]\})$
		\item $\forall I$, $\forall X$ après un $\bullet$: $J=\delta(I,X)$, numéroter si nouveau
		\item $[A\rightarrow\alpha\bullet a\beta,c]$: \textsc{Action}$[I,a]=$ shift $\delta(I,a)$
		\item $[A\rightarrow\alpha\bullet,c]$: \textsc{Action}$[I,c]=$ reduce $A\rightarrow\alpha$
		\item $[G\rightarrow S\bullet,\$]$: \textsc{Action}$[I,\$]=$ accept
		\item NT $X$: \textsc{Goto}$[I,X]=\delta(I,X)$
	\end{enumerate}
 
	Conflits (deux actions dans une case $\Rightarrow$ pas LR(1)):
	\begin{enumerate}[\roman*.]
		\item Shift/reduce: e.g. $S\rightarrow S+S\mid$ int, sur $+$
		\item Reduce/reduce
	\end{enumerate}
 
	\subsection{Exemple LR(1)}
	$G\rightarrow S\$$, $S\rightarrow CC$ ($r_\alpha$), $C\rightarrow cC$ ($r_\beta$), $C\rightarrow d$ ($r_\gamma$)
	\begin{enumerate}[\roman*.]
		\item $I_0=\{[G\rightarrow\bullet S\$,\$],[S\rightarrow\bullet CC,\$],$\\$[C\rightarrow\bullet cC,c/d],[C\rightarrow\bullet d,c/d]\}$
		\item $I_2=\delta(I_0,C)=\{[S\rightarrow C\bullet C,\$],[C\rightarrow\bullet cC,\$],$\\$[C\rightarrow\bullet d,\$]\}$
		\item $I_3=\delta(I_0,c)=\{[C\rightarrow c\bullet C,c/d],[C\rightarrow\bullet cC,c/d],$\\$[C\rightarrow\bullet d,c/d]\}$
	\end{enumerate}
  	\begin{tabular}{|c|c|c|c|c|c||c|c|c|c|c|}
		\hline
		 & $c$ & $d$ & $\$$ & $S$ & $C$ &  & $c$ & $d$ & $\$$ & $C$ \\
		\hline
		0 & s3         & s4         &     & 1 & 2 & 5 &            &            & $r_\alpha$ &   \\
		1 &            &            & acc &   &   & 6 & s6         & s7         &            & 9 \\
		2 & s6         & s7         &     &   & 5 & 7 &            &            & $r_\gamma$ &   \\
		3 & s3         & s4         &     &   & 8 & 8 & $r_\beta$  & $r_\beta$  &            &   \\
		4 & $r_\gamma$ & $r_\gamma$ &     &   &   & 9 &            &            & $r_\beta$  &   \\
		\hline
	\end{tabular}
 
	Trace \lstinline|cdd|: \lstinline|0| $\xrightarrow{s}$ \lstinline|03| $\xrightarrow{s}$ \lstinline|034| $\xrightarrow{r_\gamma}$ \lstinline|038| $\xrightarrow{r_\beta}$ \lstinline|02| $\xrightarrow{s}$ \lstinline|027| $\xrightarrow{r_\gamma}$ \lstinline|025| $\xrightarrow{r_\alpha}$ \lstinline|01| acc
  \vspace{-2em}
  \colbreak
  \section{Logique combinatoire}
	Fonction logique $f(\vec{i})=\vec{o}$; porte: son implémentation
 
	\begin{tabular}{|c|c|c|c|c|c|c|}
		\hline
		$a$ & $b$ & \textsc{and} & \textsc{or} & \textsc{xor} & \textsc{nand} & $\lnot a$ \\
		\hline
		0   & 0   & 0            & 0           & 0            & 1             & 1         \\
		0   & 1   & 0            & 1           & 1            & 1             & 1         \\
		1   & 0   & 0            & 1           & 1            & 1             & 0         \\
		1   & 1   & 1            & 1           & 0            & 0             & 0         \\
		\hline
	\end{tabular}
 
	Tout avec NAND:
	\begin{enumerate}[\roman*.]
		\item $\lnot a=\textsc{nand}(a,a)$
		\item $a\land b=\lnot\textsc{nand}(a,b)$
		\item $a\lor b=\textsc{nand}(\lnot a,\lnot b)$
		\item $a\oplus b=\textsc{nand}(\textsc{nand}(a,n),\textsc{nand}(b,n))$,\\avec $n=\textsc{nand}(a,b)$
		\item Mux $o=s\ ?\ i_1:i_0$\\$=\textsc{nand}(\textsc{nand}(\lnot s,i_0),\textsc{nand}(s,i_1))$
	\end{enumerate}
 
	Additionneurs:
	\begin{enumerate}[\roman*.]
		\item Half-adder: $s = a\oplus b$, $r = a\land b$
		\item Full-adder: $s = a\oplus b\oplus c$,\\$r = (a\land b)\lor(c\land(a\oplus b))$ ($=1$ ssi $a+b+c\geq2$)
	\end{enumerate}
 
	Bus ($k$ bits):
	\begin{enumerate}[\roman*.]
		\item Petit-boutien: $v(\vec{b})=\sum_ib_i2^i$; signe $=b_{k-1}$
		\item Addition (mod $2^k$): HA sur le bit 0, puis FA en chaîne; mux de bus: un mux par bit
	\end{enumerate}
 
	\subsection{ALU (16 bits)}
	Tout est calculé, un arbre de mux sur $op_2op_1op_0$ choisit:

	\begin{tabular}{|c|l|l|}
		\hline
		\textbf{op} & $\vec{z}$      & \textbf{Réalisation}    \\
		\hline
		0           & $x+y$          & additionneur            \\
		1           & $\sim x$       & NOT bit à bit           \\
		2           & $x\ \&\ y$     & AND bit à bit           \\
		3           & $x+1$          & retenue entrante 1      \\
		4           & $x-y$          & $x+\sim y+1$            \\
		5           & $x\oplus y$    & XOR bit à bit           \\
		6           & $x\mid y$      & OR bit à bit            \\
		7           & $x-1$          & $x+$\texttt{0xFFFF}     \\
		\hline
	\end{tabular}
 
	\colbreak
	\section{Logique séquentielle}
	Horloge $\Rightarrow$ instants $t$; état $\vec{u}_t$ ($\vec{u}_0=\vec{0}$): $(\vec{u}_{t+1},\vec{o}_t)=f(\vec{u}_t,\vec{i}_t)$.\\
	Pas une fonction mathématique; les tables ne déterminent pas le circuit.
 
	Composants:
	\begin{enumerate}[\roman*.]
		\item Registre (delay): $o_t=i_{t-1}$, $o_0=0$
		\item Mémoire 1 bit: $r_t=w_{t-1}\ ?\ v_{t-1}:r_{t-1}$\\$\Rightarrow r=\textsc{delay}(\textsc{mux}(w,r,v))$
		\item Registre 16 bits: 16 mémoires 1 bit, même $w$
	\end{enumerate}
 
	\subsection{Mémoire 2 lectures / 1 écriture}
	Lire $\vec{AL}_i$ ($k$ bits) $\rightarrow\vec{VL}_i$ (16 bits); écrire $\vec{VE}$ en $\vec{AE}$ si $e$.
 
	Récurrence sur $k$:
	\begin{enumerate}[\roman*.]
		\item $k=0$: un registre; $\vec{VL}_1=\vec{VL}_2=$ contenu
		\item $k>0$: $m_0,m_1$ sur $k-1$ bits (sans le bit $a$)
		\begin{itemize}
			\item écrire $m_b$ ssi $e\land(AE_a=b)$
			\item $\vec{VL}_i=\textsc{mux}(AL_{i,a},\ m_0.\vec{VL}_i,\ m_1.\vec{VL}_i)$
		\end{itemize}
	\end{enumerate}
 
	RAM (init. à 0) / ROM (case $i$ = tableau$[i]$) / Mixte (ROM programme + RAM données).
 
	\colbreak
	\section{CPU}
	Composants:
	\begin{enumerate}[\roman*.]
		\item Mémoire RAM/ROM: $2^{12}$ mots de 16 bits
		\item $2^4=16$ registres de 16 bits; \lstinline|PC|; terminal $80\times24$
	\end{enumerate}
 
	Instruction (16 bits): \lstinline|OP + 16*(R1 + 16*(R2 + 16*R3))|
	\begin{enumerate}[\roman*.]
		\item Bits 0--3: OP / 4--7: $R_1$ (écrit)
		\item Bits 8--11: $R_2$ / 12--15: $R_3$ (lus)
	\end{enumerate}
 
	\begin{tabular}{|c|l|l|}
		\hline
		\textbf{OP} & \textbf{Nom} & \textbf{Effet}                              \\
		\hline
		$\geq8$     & ALU          & $R[R_1]=ALU(OP{-}8,R[R_2],R[R_3])$          \\
		7           & CST HI       & $R[R_1]=\overline{R_3R_2}\times2^8$         \\
		6           & CST LOW      & $R[R_1]=\overline{R_3R_2}$                  \\
		5           & LOAD         & $R[R_1]=RAM[R[R_2]]$                        \\
		4           & PC           & $R[R_1]=PC$                                 \\
		3           & OUTPUT       & $Out(R[R_2],R[R_3])$                        \\
		2           & STORE        & $RAM[R[R_2]]=R[R_3]$                        \\
		1           & JL           & $PC=R[R_3]$ si $(R[R_2]<0)\oplus R_1[0]$    \\
		0           & JZ           & $PC=R[R_3]$ si $(R[R_2]=0)\oplus R_1[0]$    \\
		\hline
	\end{tabular}
	\begin{enumerate}[\roman*.]
		\item $\overline{R_3R_2}=R_2+16R_3$ (8 bits)
		\item $Out(x,y)$: caractère $y$ en $(x\%80,\ x/80)$
		\item ALU: add 8, inv 9, et 10, inc 11, sub 12, xor 13, ou 14, dec 15
		\item $R_1[0]=1$: saut inversé (non nul / non négatif)
	\end{enumerate}
 
	\subsection{Un cycle}
	\begin{enumerate}[\roman*.]
		\item Fetch: $ins=MEM[PC]$; decode OP, $R_1$, $R_2$, $R_3$
		\item Lire $A=R[R_2]$, $B=R[R_3]$ (toujours, sans risque)
		\item Calculer $ALU(OP_{0..2},A,B)$ et $D=MEM[A]$
		\item Écrire $R[R_1]$ si $OP\geq4$; $MEM[A]=B$ si $OP=2$; sortie si $OP=3$
		\item $PC\leftarrow saut\ ?\ B:PC+1$
	\end{enumerate}
 
	Exemple (affiche A en boucle):\\
	\noindent\begin{minipage}{\linewidth}
		\begin{lstlisting}[style=ocaml]
ins_cst_low 0 0 ;    (* r0 = 0 *)
ins_cst_low (int_of_char 'A') 1 ;
ins_output 0 1 ;     (* Out(r0, r1) *)
ins_cst_low 2 2 ;    (* r2 = 2 *)
ins_jump_zero 0 2    (* PC = 2 *)
\end{lstlisting}
	\end{minipage}
 
	%%%%%%%%%%%%%%%%%%%%%%%%%%%%%%%%%%%%%%%%%%%%%%%%%%%%%%
	%                       End                          %
	%%%%%%%%%%%%%%%%%%%%%%%%%%%%%%%%%%%%%%%%%%%%%%%%%%%%%%
\end{multicols*}
\end{document}
